\section{Discussion}
\label{sec:discussion}

\subsection{The Coupling Disparity Resolves a Paradox}

The central result of this study is the $10^3$--$10^4$ factor separating
mechanical from airborne coupling efficiency.  This disparity has a simple
physical origin: mechanical vibration drives the shell wall directly, while
airborne sound must couple through the acoustic radiation impedance mismatch.
In the long-wavelength limit ($ka \ll 1$), this mismatch imposes a $(ka)^n$
penalty on flexural mode $n$, which at \SI{7}{\hertz} amounts to a factor of
$2.8 \times 10^{-4}$ for $n = 2$.

This finding resolves an apparent paradox in the literature: whole-body
vibration at \SIrange{4}{8}{\hertz} is well documented to cause
gastrointestinal distress (including nausea, abdominal pain, and in extreme
cases, incontinence) \cite{griffin1990handbook,mansfield2005human}, whereas airborne infrasound studies at levels up to
\SI{150}{\decibel} report no comparable effects \cite{Leventhall2007}.The modal analysis shows that
both observations are consistent with the same underlying resonance---but the
two excitation pathways differ by three to four orders of magnitude in
coupling efficiency.

\subsection{Implications for the ``Brown Note'' Hypothesis}

The results indicate that abdominal resonance in the \SIrange{4}{10}{\hertz}
range is physically real and well-supported by ISO~2631 data.  The ``brown
note'' hypothesis is therefore half-correct: the resonance exists.  However,
\emph{airborne} excitation of this resonance is energetically negligible at
any plausible SPL.  The effects attributed to infrasound are instead explained
by mechanical vibration, which is the actual exposure pathway in the
occupational scenarios (heavy equipment, transport vehicles) where
gastrointestinal symptoms are reported.

\subsection{Gas Pocket Mechanism}

We note that bowel gas pockets (radius $R \approx \SI{2}{\centi\metre}$) have
Minnaert resonance frequencies in the range \SIrange{30}{330}{\hertz}.  At
infrasonic frequencies, sub-resonant but non-negligible acoustic displacements
of order \SI{1}{\micro\metre} at \SI{120}{\decibel} may occur within these
pockets.  This mechanism could potentially explain inter-individual
variability in infrasound sensitivity, but the present analysis cannot make a
strong quantitative claim without modelling the gas--tissue interaction in
detail.

\subsection{Limitations}

\begin{enumerate}
  \item \textbf{Boundary conditions}: The free-shell model provides a lower
    bound on mode frequencies. Partial clamping by the spine and pelvis will
    raise the $n = 2$ frequency, potentially by a factor of 2--3$\times$.
    Finite element validation is needed to resolve this uncertainty.
  \item \textbf{Linear theory}: The small-displacement assumption may not hold
    for large-amplitude mechanical vibration.
  \item \textbf{Homogeneous wall}: The abdominal wall is a multi-layer
    composite, though our laminate analysis shows this introduces less than 5\%
    error in the eigenfrequency.
  \item \textbf{Thin-shell theory}: The wall-thickness-to-radius ratio
    ($h/R \approx 0.06$) is near the upper bound of thin-shell validity.
  \item \textbf{No organ-level model}: Internal organs are treated as fluid;
    a more realistic model would include solid inclusions.
\end{enumerate}
